% ==========================================================
% 100_production_implications.tex
% Forecast Admissibility Surface (FAS)
% ==========================================================

\section{Implications for Production Management}
\label{sec:fas_production_implications}

The Forecast Admissibility Surface has direct implications for how forecasting
outputs are interpreted and applied in operational production management
settings. By explicitly distinguishing between admissible and inadmissible
slices, FAS prevents structurally unstable behavior from propagating into
planning, execution, and control decisions.

\subsection{Preventing False Precision}

In many production environments, forecasts are consumed as point estimates or
narrow ranges that implicitly convey confidence. When applied to structurally
hostile slices, such forecasts can create a false sense of precision, leading to
overconfident staffing, procurement, or preparation decisions. FAS mitigates
this risk by blocking or constraining forecasts in regimes where baseline error
anatomy indicates that precision is illusory.

\subsection{Targeted Use of Forecasting Capacity}

Forecasting capacity—both computational and human—is finite. FAS enables systems
to allocate this capacity where it is most likely to yield stable and actionable
signal. By excluding \Blocked\ slices and flagging \Conditional\ slices for
restricted handling, modeling effort can be focused on demand regimes that
benefit meaningfully from learning and evaluation.

\subsection{Operational Handling of Inadmissible Slices}

Blocking a slice from forecasting does not imply that the corresponding
operational process can be ignored. Rather, it signals that alternative handling
mechanisms are required. These may include deterministic rules, manual planning,
safety buffers, or exogenous overrides informed by domain expertise. FAS makes
explicit where such non-forecast-based strategies are necessary.

\subsection{Stability of Downstream Signals}

By constraining forecasting to structurally admissible slices, FAS improves the
stability of downstream metrics and governance signals. Evaluation diagnostics
become less sensitive to rare extremes, readiness adjustments behave more
consistently, and policy decisions are less likely to oscillate in response to
pathological behavior. This stability is particularly important in production
management contexts where frequent reversals erode trust in analytic systems.

\subsection{Interpretive Discipline}

Finally, FAS enforces interpretive discipline in production environments. It
clarifies where forecasts should be trusted, where they should be treated with
caution, and where they should not be used at all. By making these boundaries
explicit and auditable, FAS supports more informed and resilient operational
decision-making.
